\chapter{Was gewesen wäre}

\section{Der Leitfall, einzeln betrachtet}

\begin{merksatz}
\boxtitel{Ein Patient}
Herr K., $68$ Jahre, Sepsis, schwer krank. Er erhielt das Antibiotikum nach vierzig Minuten und überlebte.

\begin{center}
\begin{minipage}{0.9\textwidth}
\textbf{Hat die frühe Gabe ihm geholfen?}
\end{minipage}
\end{center}

\begin{center}
\begin{minipage}{0.88\textwidth}
\kursiv{Die Frage ist klar gestellt und dennoch nicht zu beantworten. \textbf{Um sie zu beantworten, müsste man wissen, was mit ihm geschehen wäre, hätte man länger gewartet} --- und das kann niemand wissen, weil es diesen Patienten nur einmal gibt. Dieses Kapitel gibt der unbeobachtbaren Größe einen Namen, zeigt, warum sie für den Einzelnen unerreichbar bleibt, und unter welchen Bedingungen sie für eine Gruppe dennoch schätzbar ist.}
\end{minipage}
\end{center}
\end{merksatz}

\begin{definition}
\boxtitelm{Definition 2.1 --- potentielle Ergebnisse}
Für jeden Patienten $i$ und jede Behandlung $x$ sei
\[ Y_i(x) \]
das Ergebnis, das eintreten \kursiv{würde}, wenn er die Behandlung $x$ erhielte.

Bei zwei Behandlungen hat jeder Patient zwei solche Größen, $Y_i(1)$ und $Y_i(0)$, und der \Begriff{individuelle Effekt} ist
\[ \tau_i = Y_i(1) - Y_i(0) . \]
\end{definition}

\begin{merksatz}
\boxtitel{Was daran neu ist}
\begin{center}
{\small
\renewcommand{\arraystretch}{1.55}
\begin{tabular}{@{}P{2.6cm}P{6.0cm}@{}}
\toprule
\textbf{Bisher} & Eine Zufallsgröße $Y$ nimmt Werte an, und man fragt nach ihrer Verteilung, auch bedingt auf anderes.\\
\textbf{Jetzt} & Jeder Patient trägt \kursiv{zwei} Werte in sich --- einen für jede mögliche Behandlung. Beobachtbar ist immer nur einer.\\
\bottomrule
\end{tabular}}
\end{center}

\schritt{Die Verbindung zum Beobachteten}
\[ Y_i = Y_i(X_i) \]
--- das beobachtete Ergebnis ist dasjenige potentielle Ergebnis, das zur tatsächlich gegebenen Behandlung gehört. Diese Annahme heißt \Begriff{Konsistenz}.

\begin{center}
\begin{minipage}{0.88\textwidth}
\kursiv{Sie klingt selbstverständlich und ist es nicht. \textbf{Sie verlangt, dass \enquote{die Behandlung} eindeutig bestimmt ist}: Wenn \enquote{früh} bei dem einen vierzig Minuten und bei dem anderen fünfzig heißt, so gibt es nicht ein $Y_i(1)$, sondern viele. Der Bruchfall kommt darauf zurück.}
\end{minipage}
\end{center}
\end{merksatz}

\begin{merksatz}
\boxtitel{Satz 2.2 --- das grundlegende Problem}
Der individuelle Effekt $\tau_i$ ist für keinen einzigen Patienten beobachtbar.

\schritt{Begründung}
Man beobachtet $Y_i(1)$ oder $Y_i(0)$, nie beides --- der Patient wurde entweder früh behandelt oder nicht.

\begin{center}
{\small
\renewcommand{\arraystretch}{1.5}
\begin{tabular}{@{}P{1.8cm}P{1.6cm}P{1.8cm}P{1.8cm}@{}}
\toprule
\textbf{Patient} & $X_i$ & $Y_i(1)$ & $Y_i(0)$\\
\midrule
Herr K. & $1$ & \textcolor{gruen}{überlebt} & \textcolor{akzent}{?}\\
Frau L. & $0$ & \textcolor{akzent}{?} & \textcolor{gruen}{gestorben}\\
\bottomrule
\end{tabular}}
\end{center}

\begin{center}
\begin{minipage}{0.9\textwidth}
\textbf{Die Hälfte der Tabelle fehlt --- und zwar unvermeidlich, nicht aus Nachlässigkeit.}
\end{minipage}
\end{center}

\begin{center}
\begin{minipage}{0.88\textwidth}
\kursiv{\Person{Paul Holland} nannte das 1986 das \enquote{grundlegende Problem der kausalen Inferenz}. \textbf{Es ist kein Messproblem, sondern ein begriffliches}: Die fehlende Größe existiert nicht als Tatsache über die Welt, sondern nur als Aussage darüber, was gewesen wäre. Man kann sie nicht nachmessen, auch nicht mit beliebigem Aufwand.}
\end{minipage}
\end{center}
\end{merksatz}

\section{Vom Einzelnen zur Gruppe}

\begin{definition}
\boxtitelm{Definition 2.3 --- Durchschnittseffekte}
\begin{center}
{\small
\renewcommand{\arraystretch}{1.6}
\begin{tabular}{@{}P{2.8cm}P{5.8cm}@{}}
\toprule
$\tau = \mathbb{E}[Y(1)-Y(0)]$ & \textbf{Durchschnittseffekt} über alle.\\
$\tau_{\text{beh}} = \mathbb{E}[Y(1)-Y(0) \mid X=1]$ & \textbf{Effekt bei den Behandelten}.\\
\bottomrule
\end{tabular}}
\end{center}
\end{definition}

\begin{merksatz}
\boxtitel{Warum zwei Größen}
\begin{center}
{\small
\renewcommand{\arraystretch}{1.5}
\begin{tabular}{@{}P{2.6cm}P{6.0cm}@{}}
\toprule
$\tau$ & Was geschähe, wenn man \kursiv{alle} behandelte? Die Frage einer Leitlinie.\\
$\tau_{\text{beh}}$ & Was hat die Behandlung denen gebracht, die sie erhalten haben? Die Frage einer Nachbetrachtung.\\
\bottomrule
\end{tabular}}
\end{center}

\begin{center}
\begin{minipage}{0.88\textwidth}
\kursiv{Beide fallen auseinander, sobald die Behandelten sich von den übrigen unterscheiden --- also fast immer. \textbf{Im Leitfall sind die Behandelten überwiegend schwer krank}; dort kann die Wirkung größer oder kleiner sein als bei den Leichtkranken. Aufgabe 2.1 rechnet beide aus und findet $-0{,}25$ gegen $-0{,}375$.}
\end{minipage}
\end{center}
\end{merksatz}

\section{Vier Annahmen}

\begin{merksatz}
\boxtitel{Was man braucht, um zu rechnen}
\begin{center}
{\small
\renewcommand{\arraystretch}{1.6}
\begin{tabular}{@{}P{2.8cm}P{5.8cm}@{}}
\toprule
\textbf{Konsistenz} & $Y_i = Y_i(X_i)$ --- die Behandlung ist eindeutig bestimmt.\\
\textbf{Keine Wechselwirkung} & Die Behandlung des einen wirkt nicht auf das Ergebnis des anderen.\\
\textbf{Positivität} & In jeder Gruppe kommen beide Behandlungen vor: $0 < \Pr[X=1\mid Z] < 1$.\\
\textbf{Bedingte Unabhängigkeit} & $\bigl(Y(1),Y(0)\bigr) \perp X \mid Z$ --- bei gleichem $Z$ ist die Behandlung so gut wie ausgelost.\\
\bottomrule
\end{tabular}}
\end{center}

\begin{center}
\begin{minipage}{0.9\textwidth}
\textbf{Nur die dritte Annahme lässt sich an den Daten prüfen.}
\end{minipage}
\end{center}

\begin{center}
\begin{minipage}{0.88\textwidth}
\kursiv{Die erste ist eine Frage der Begriffsbildung, die zweite eine des Sachverhalts, die vierte eine des Fachwissens --- und gerade die vierte trägt die ganze Last. \textbf{Sie besagt, dass es außer $Z$ keine gemeinsame Ursache von Behandlung und Ergebnis gibt}, und das ist genau die unprüfbare Annahme aus dem Bruchfall von Kapitel~1. Kapitel~6 gibt ihr eine grafische Gestalt, an der man sie wenigstens diskutieren kann.}
\end{minipage}
\end{center}
\end{merksatz}

\begin{merksatz}
\boxtitel{Satz 2.4 --- die Anpassungsformel}
Unter Konsistenz, Positivität und bedingter Unabhängigkeit gilt
\[ \mathbb{E}\bigl[Y(x)\bigr] = \sum_z \mathbb{E}\bigl[Y \mid X=x,\ Z=z\bigr]\cdot\Pr[Z=z] . \]

\schritt{Beweis}
\kursiv{Erstens --- nach $Z$ zerlegen} (Satz von der totalen Erwartung, Band 4, Kapitel~9):
\[ \mathbb{E}[Y(x)] = \sum_z \mathbb{E}[Y(x) \mid Z=z]\,\Pr[Z=z] \]

\kursiv{Zweitens --- die Behandlung hinzunehmen:} Nach der bedingten Unabhängigkeit ändert die Bedingung auf $X=x$ bei gegebenem $Z=z$ nichts an der Verteilung von $Y(x)$:
\[ \mathbb{E}[Y(x) \mid Z=z] = \mathbb{E}[Y(x) \mid X=x,\ Z=z] \]
Damit dieser Ausdruck überhaupt definiert ist, braucht man die Positivität.

\kursiv{Drittens --- Konsistenz:} Unter denen mit $X=x$ ist $Y(x)$ das beobachtete $Y$:
\[ \mathbb{E}[Y(x) \mid X=x,\ Z=z] = \mathbb{E}[Y \mid X=x,\ Z=z] \]

\kursiv{Einsetzen} liefert die Behauptung. $\blacksquare$

\begin{center}
\begin{minipage}{0.9\textwidth}
\textbf{Jede der drei Annahmen wird in genau einem Schritt gebraucht --- und ohne sie bricht der Beweis an dieser Stelle ab.}
\end{minipage}
\end{center}

\begin{center}
\begin{minipage}{0.88\textwidth}
\kursiv{\textbf{Das ist die Standardisierung aus Kapitel~1}, nun mit Begründung. Sie war dort eine plausible Rechnung; hier ist sie ein Satz mit angebbaren Voraussetzungen. Der Unterschied zwischen beidem ist der Gegenstand dieses Bandes.}
\end{minipage}
\end{center}
\end{merksatz}

\section{Der Bruchfall}

\begin{warnung}
\boxtitelr{Wechselwirkung: der Fall, in dem alles zusammenbricht}
Die zweite Annahme --- die Behandlung des einen wirke nicht auf das Ergebnis des anderen --- ist bei ansteckenden Krankheiten falsch.

\begin{center}
\begin{minipage}{0.88\textwidth}
\kursiv{Wird ein Patient geimpft, so sinkt auch das Risiko seiner Nachbarn. \textbf{Dann gibt es nicht zwei potentielle Ergebnisse je Person, sondern eines für jede denkbare Belegung der ganzen Bevölkerung} --- bei $n$ Personen also $2^{n}$. Definition 2.1 ist dann gar nicht anwendbar.}
\end{minipage}
\end{center}

\schritt{Was man stattdessen tut}
\begin{center}
{\small
\renewcommand{\arraystretch}{1.5}
\begin{tabular}{@{}P{2.6cm}P{6.0cm}@{}}
\toprule
\textbf{Gruppen bilden} & Man randomisiert ganze Dörfer oder Stationen statt einzelner Personen.\\
\textbf{Effekte trennen} & Direkter Effekt der eigenen Impfung, indirekter Effekt der Impfung anderer.\\
\bottomrule
\end{tabular}}
\end{center}

\begin{center}
\begin{minipage}{0.88\textwidth}
\kursiv{\textbf{Der indirekte Effekt ist die Herdenwirkung aus Band 18, Kapitel~7.} Dort wurde sie als Eigenschaft des SIR-Modells hergeleitet; hier erscheint sie als Grund, warum ein ganzes Begriffsgebäude nicht greift. Bei Antibiotika kommt eine zweite Wechselwirkung hinzu: Die Resistenzlage einer Station hängt davon ab, wie andere Patienten behandelt wurden --- Band 21, Kapitel~10.}
\end{minipage}
\end{center}
\end{warnung}

\begin{warnung}
\boxtitelr{Konsistenz: was heißt \enquote{früh}?}
\begin{center}
{\small
\renewcommand{\arraystretch}{1.52}
\begin{tabular}{@{}P{2.6cm}P{6.0cm}@{}}
\toprule
\textbf{Unscharf} & \enquote{Innerhalb einer Stunde} umfasst fünf Minuten und neunundfünfzig.\\
\textbf{Mit Begleitumständen} & Wer früh behandelt wird, bekommt oft auch früher Flüssigkeit und mehr Aufmerksamkeit.\\
\bottomrule
\end{tabular}}
\end{center}

\begin{center}
\begin{minipage}{0.88\textwidth}
\kursiv{Der geschätzte Effekt gilt dann für das Bündel, nicht für das Antibiotikum allein. \textbf{Die Regel lautet: Man muss den Eingriff so genau beschreiben können, dass zwei Ärzte ihn gleich ausführen würden.} Bei Medikamenten gelingt das meist, bei Begriffen wie \enquote{Adipositas} oder \enquote{Stress} nicht --- und deshalb sind kausale Aussagen über solche Größen schwer zu fassen.}
\end{minipage}
\end{center}
\end{warnung}

\section{Übungsaufgaben}

\begin{aufgabe}
\boxtitel{Aufgabe 2.1 --- Die vollständige Tabelle}
Für acht Patienten seien ausnahmsweise \kursiv{beide} potentiellen Ergebnisse bekannt ($1$ bedeutet gestorben):
\begin{center}
{\small
\renewcommand{\arraystretch}{1.44}
\begin{tabular}{@{}P{1.2cm}P{1.8cm}P{1.2cm}P{1.2cm}P{1.0cm}P{1.4cm}@{}}
\toprule
\textbf{Nr.} & \textbf{Schwere} & $Y(0)$ & $Y(1)$ & $X$ & \textbf{$Y$ beob.}\\
\midrule
$1$ & schwer & $1$ & $1$ & $1$ & $1$\\
$2$ & schwer & $1$ & $0$ & $1$ & $0$\\
$3$ & schwer & $1$ & $1$ & $1$ & $1$\\
$4$ & schwer & $1$ & $0$ & $0$ & $1$\\
$5$ & leicht & $0$ & $0$ & $1$ & $0$\\
$6$ & leicht & $1$ & $0$ & $0$ & $1$\\
$7$ & leicht & $0$ & $0$ & $0$ & $0$\\
$8$ & leicht & $0$ & $0$ & $0$ & $0$\\
\bottomrule
\end{tabular}}
\end{center}
Berechnen Sie $\tau$, $\tau_{\text{beh}}$, den naiven Vergleich und die Anpassungsformel.

\loesung
\schritt{Schritt 1 --- Die individuellen Effekte}
\begin{center}
{\small
\renewcommand{\arraystretch}{1.44}
\begin{tabular}{@{}P{1.2cm}P{1.2cm}P{1.2cm}P{1.2cm}P{1.2cm}P{1.2cm}P{1.2cm}P{1.2cm}@{}}
\toprule
$1$ & $2$ & $3$ & $4$ & $5$ & $6$ & $7$ & $8$\\
\midrule
$0$ & $-1$ & $0$ & $-1$ & $0$ & $-1$ & $0$ & $0$\\
\bottomrule
\end{tabular}}
\end{center}

\begin{center}
\begin{minipage}{0.9\textwidth}
\textbf{Bei drei von acht Patienten rettet die Behandlung das Leben, bei fünf ändert sie nichts. Sie schadet niemandem.}
\end{minipage}
\end{center}

\schritt{Schritt 2 --- Die Durchschnittseffekte}
\[ \tau = \frac{-3}{8} = -0{,}375 \]
\kursiv{Behandelt sind $1$, $2$, $3$, $5$:}
\[ \tau_{\text{beh}} = \frac{0-1+0+0}{4} = -0{,}250 \]
\kursiv{Unbehandelt sind $4$, $6$, $7$, $8$:}
\[ \tau_{\text{unbeh}} = \frac{-1-1+0+0}{4} = -0{,}500 \]

\begin{center}
\begin{minipage}{0.88\textwidth}
\kursiv{\textbf{Die Behandlung hätte bei den Unbehandelten doppelt so viel genützt wie bei den Behandelten.} Wer nur $\tau_{\text{beh}}$ berichtet, unterschätzt, was eine allgemeine Leitlinie brächte.}
\end{minipage}
\end{center}

\schritt{Schritt 3 --- Der naive Vergleich}
\[ \bar Y_{X=1} = \frac{1+0+1+0}{4} = 0{,}50, \qquad \bar Y_{X=0} = \frac{1+1+0+0}{4} = 0{,}50 \]
\[ \text{Unterschied} = 0 \]

\begin{center}
\begin{minipage}{0.9\textwidth}
\textbf{Der naive Vergleich zeigt gar keinen Effekt --- der wahre ist $-0{,}375$.}
\end{minipage}
\end{center}

\schritt{Schritt 4 --- Die Anpassungsformel}
\begin{center}
{\small
\renewcommand{\arraystretch}{1.46}
\begin{tabular}{@{}P{1.8cm}P{2.4cm}P{2.4cm}P{1.6cm}@{}}
\toprule
\textbf{Gruppe} & \textbf{behandelt} & \textbf{unbehandelt} & \textbf{Anteil}\\
\midrule
schwer & $2/3 = 0{,}667$ & $1/1 = 1{,}000$ & $0{,}5$\\
leicht & $0/1 = 0{,}000$ & $1/3 = 0{,}333$ & $0{,}5$\\
\bottomrule
\end{tabular}}
\end{center}

\[ \hat\tau = 0{,}5\,(0{,}667-1{,}000) + 0{,}5\,(0{,}000-0{,}333) = -0{,}167 - 0{,}167 = -0{,}333 \]

\begin{center}
\begin{minipage}{0.88\textwidth}
\kursiv{\textbf{Die Anpassung liefert $-0{,}333$ gegen einen wahren Wert von $-0{,}375$} --- der Rest ist Zufall der kleinen Stichprobe, nicht Verzerrung: In jeder Gruppe stehen nur ein bis drei Patienten. Bei $1200$ Patienten wie in Kapitel~1 verschwindet dieser Unterschied.}
\end{minipage}
\end{center}
\end{aufgabe}

\begin{aufgabe}
\boxtitel{Aufgabe 2.2 --- Die Annahme prüfen}
Zeigen Sie an derselben Tabelle, dass die bedingte Unabhängigkeit hier tatsächlich gilt, und was geschähe, wenn sie verletzt wäre.

\loesung
\schritt{Schritt 1 --- Was zu prüfen ist}
Innerhalb jeder Schweregruppe müssten Behandelte und Unbehandelte dieselben potentiellen Ergebnisse haben.

\begin{center}
{\small
\renewcommand{\arraystretch}{1.46}
\begin{tabular}{@{}P{1.8cm}P{2.6cm}P{2.6cm}@{}}
\toprule
\textbf{Gruppe} & $\bar Y(0)$ behandelt & $\bar Y(0)$ unbehandelt\\
\midrule
schwer & $1{,}000$ & $1{,}000$\\
leicht & $0{,}000$ & $0{,}333$\\
\bottomrule
\end{tabular}}
\end{center}

\begin{center}
\begin{minipage}{0.9\textwidth}
\textbf{In der schweren Gruppe stimmen sie überein, in der leichten nicht ganz} --- $0{,}000$ gegen $0{,}333$.
\end{minipage}
\end{center}

\schritt{Schritt 2 --- Ist das eine Verletzung?}
\begin{center}
\begin{minipage}{0.88\textwidth}
\kursiv{Nicht notwendig. Die leichte Gruppe umfasst vier Patienten, davon einen behandelten. \textbf{Ein Unterschied von einem Drittel bei so wenigen Fällen ist mit Zufall vereinbar.} Die Annahme ist eine Aussage über die Verteilung, nicht über die Stichprobe --- und genau deshalb lässt sie sich an der Stichprobe nicht widerlegen.}
\end{minipage}
\end{center}

\schritt{Schritt 3 --- Eine echte Verletzung}
Man tausche die Behandlung der Patienten $5$ und $6$: Nun ist Patient $6$ behandelt (er hätte ohne Behandlung $Y(0)=1$) und Patient $5$ nicht.

\begin{center}
{\small
\renewcommand{\arraystretch}{1.46}
\begin{tabular}{@{}P{1.8cm}P{2.4cm}P{2.4cm}P{1.6cm}@{}}
\toprule
\textbf{Gruppe} & \textbf{behandelt} & \textbf{unbehandelt} & \textbf{Anteil}\\
\midrule
schwer & $0{,}667$ & $1{,}000$ & $0{,}5$\\
leicht & $0/1 = 0{,}000$ & $0/3 = 0{,}000$ & $0{,}5$\\
\bottomrule
\end{tabular}}
\end{center}

\[ \hat\tau = 0{,}5\,(-0{,}333) + 0{,}5\,(0{,}000) = -0{,}167 \]

\begin{center}
\begin{minipage}{0.9\textwidth}
\textbf{Derselbe wahre Effekt von $-0{,}375$, nun auf $-0{,}167$ geschätzt} --- weil innerhalb der leichten Gruppe gerade der Kränkere behandelt wurde.
\end{minipage}
\end{center}

\schritt{Schritt 4 --- Die Lehre}
\begin{center}
\begin{minipage}{0.88\textwidth}
\kursiv{\textbf{Die Anpassung an $Z$ hilft nur, wenn $Z$ alles enthält, wonach ausgewählt wurde.} Hier hat der Arzt innerhalb der leichten Gruppe nach etwas ausgewählt, das nicht in $Z$ steht --- vielleicht dem Aussehen des Patienten. Kein Rechenverfahren kann das bemerken; nur die Frage, woraufhin der Arzt entschieden hat.}
\end{minipage}
\end{center}
\end{aufgabe}

\begin{aufgabe}
\boxtitel{Aufgabe 2.3 --- Positivität}
In einer Klinik erhalten alle Schwerkranken das Antibiotikum früh. Was folgt daraus, und was geschieht bei $98$ statt $100$ Prozent?

\loesung
\schritt{Schritt 1 --- Bei hundert Prozent}
In der Anpassungsformel steht für die schwere Gruppe der Ausdruck
\[ \mathbb{E}[Y \mid X=0,\ Z=\text{schwer}] \]
--- ein Mittelwert über eine leere Menge.

\begin{center}
\begin{minipage}{0.9\textwidth}
\textbf{Der Durchschnittseffekt ist nicht schätzbar.} Nicht ungenau, sondern gar nicht.
\end{minipage}
\end{center}

\schritt{Schritt 2 --- Was trotzdem geht}
\begin{center}
{\small
\renewcommand{\arraystretch}{1.5}
\begin{tabular}{@{}P{2.8cm}P{5.8cm}@{}}
\toprule
\textbf{Andere Zielgruppe} & Der Effekt bei den Leichtkranken bleibt schätzbar.\\
\textbf{Modellannahme} & Man unterstellt, der Effekt sei in beiden Gruppen gleich, und überträgt. Das ist eine Annahme, keine Schätzung.\\
\bottomrule
\end{tabular}}
\end{center}

\schritt{Schritt 3 --- Bei achtundneunzig Prozent}
Von $600$ Schwerkranken sind $12$ unbehandelt. Der Mittelwert über diese $12$ hat die Streuung
\[ \frac{\sigma}{\sqrt{12}} = \frac{0{,}5}{3{,}46} = 0{,}144 \]
bei einem Anteil nahe $0{,}5$.

\begin{center}
\begin{minipage}{0.9\textwidth}
\textbf{Ein Vertrauensintervall von etwa $\pm28$ Prozentpunkten --- für eine Gruppe, die die Hälfte des Gewichts trägt.}
\end{minipage}
\end{center}

\schritt{Schritt 4 --- Die Verfahren aus Kapitel 8}
\begin{center}
\begin{minipage}{0.88\textwidth}
\kursiv{Gewichtet man Patienten mit dem Kehrwert ihrer Behandlungswahrscheinlichkeit --- Kapitel~8 ---, so erhält jeder der zwölf das Gewicht $1/0{,}02 = 50$. \textbf{Zwölf Patienten bestimmen dann die Hälfte des Ergebnisses}, und ein einzelner Todesfall verschiebt es um mehrere Prozentpunkte.}
\end{minipage}
\end{center}

\begin{center}
{\small
\renewcommand{\arraystretch}{1.48}
\begin{tabular}{@{}P{2.4cm}P{2.4cm}@{}}
\toprule
\textbf{Anteil behandelt} & \textbf{Gewicht}\\
\midrule
$50\,\%$ & $2$\\
$90\,\%$ & $10$\\
$98\,\%$ & $50$\\
\bottomrule
\end{tabular}}
\end{center}

\begin{center}
\begin{minipage}{0.88\textwidth}
\kursiv{\textbf{Die Positivität ist keine Ja-Nein-Eigenschaft, sondern eine Frage des Grades.} Je näher man ihrer Grenze kommt, desto mehr hängt das Ergebnis an wenigen Patienten --- ein Übergang von Verzerrung zu Varianz, wie in Band 22, Kapitel~2. Die übliche Abhilfe ist, die Zielgruppe zu beschränken und den Effekt nur dort zu berichten, wo beide Behandlungen tatsächlich vorkommen.}
\end{minipage}
\end{center}
\end{aufgabe}

\begin{geschichte}
\boxtitelg{Eine Notation aus der Landwirtschaft}
\Person{Jerzy Neyman} führte die potentiellen Ergebnisse 1923 in seiner Doktorarbeit ein --- für Feldversuche mit Kartoffelsorten.

\begin{center}
\begin{minipage}{0.9\textwidth}
\kursiv{Er schrieb den Ertrag jeder Parzelle für jede mögliche Sorte an und behandelte die Zuteilung als Ziehung ohne Zurücklegen. \textbf{Die Arbeit erschien auf Polnisch und blieb sechzig Jahre unbeachtet}; erst 1990 wurde sie ins Englische übersetzt.}
\end{minipage}
\end{center}

\schritt{Rubin}
\Person{Donald Rubin} griff die Notation 1974 auf und übertrug sie auf Beobachtungsdaten.

\begin{center}
\begin{minipage}{0.9\textwidth}
\kursiv{Sein Beitrag war die Einsicht, dass man dieselbe Sprache benutzen kann, wenn nicht das Los, sondern ein Arzt entscheidet --- sofern man angeben kann, woraufhin er entschieden hat. \textbf{Das Gebäude heißt heute nach Neyman und Rubin}, und die vier Annahmen dieses Kapitels stammen im Wesentlichen aus seinen Arbeiten der 1970er und 1980er Jahre.}
\end{minipage}
\end{center}

\schritt{Holland}
\begin{center}
\begin{minipage}{0.9\textwidth}
\kursiv{\Person{Paul Holland} fasste 1986 den Stand in einem Aufsatz zusammen, der den Satz enthält, \enquote{keine Ursache ohne Manipulation}. \textbf{Danach wären Aussagen wie \enquote{sie wurde wegen ihres Geschlechts nicht befördert} keine kausalen Aussagen}, weil man das Geschlecht nicht zuteilen kann. Der Satz ist seither heftig umstritten; die heutige Mehrheitsmeinung hält ihn für zu eng, verlangt aber, dass man genau angibt, welcher Eingriff gemeint ist.}
\end{minipage}
\end{center}
\end{geschichte}

\begin{forschung}
\boxtitelp{Effekte heute}
\begin{center}
{\small
\renewcommand{\arraystretch}{1.45}
\begin{tabular}{@{}P{2.8cm}P{5.4cm}@{}}
\toprule
\textbf{Schranken} & Was sich über den Effekt sagen lässt, wenn die Annahmen \kursiv{nicht} gelten.\\
\textbf{Empfindlichkeit} & Wie stark eine unbeobachtete Störgröße wirken müsste, um das Ergebnis umzukehren.\\
\textbf{Wechselwirkung} & Effekte in Netzwerken, wo die Nachbarn mitwirken.\\
\bottomrule
\end{tabular}}
\end{center}

\schritt{Zur zweiten Zeile}
\begin{center}
\begin{minipage}{0.88\textwidth}
\kursiv{Statt die bedingte Unabhängigkeit zu behaupten, fragt man: Wie stark müsste eine unbeobachtete Größe mit Behandlung \kursiv{und} Ergebnis zusammenhängen, damit der geschätzte Effekt verschwindet? \textbf{Für den Zusammenhang zwischen Rauchen und Lungenkrebs hätte eine solche Größe das Risiko verneunfachen müssen} --- ein Argument, das \Person{Jerome Cornfield} 1959 gegen Fisher vorbrachte und das heute in verfeinerter Form Standard ist.}
\end{minipage}
\end{center}
\end{forschung}

\begin{merksatz}
\boxtitel{Was Sie jetzt können}
\begin{center}
\renewcommand{\arraystretch}{1.5}
\begin{tabular}{@{}P{3.0cm}P{5.6cm}@{}}
\toprule
\textbf{Definition 2.1} & zwei Werte je Patient, einer sichtbar\\
\textbf{Satz 2.2} & der individuelle Effekt bleibt unbekannt\\
\textbf{Definition 2.3} & $\tau$ und $\tau_{\text{beh}}$ fallen auseinander\\
\textbf{Vier Annahmen} & nur die Positivität ist prüfbar\\
\textbf{Satz 2.4} & Anpassungsformel, dreimal begründet\\
\bottomrule
\end{tabular}
\end{center}

\begin{center}
\begin{minipage}{0.88\textwidth}
\kursiv{Die schwierigste Annahme war die bedingte Unabhängigkeit. Das nächste Kapitel zeigt den einen Fall, in dem sie nicht behauptet, sondern hergestellt wird --- durch das Los.}
\end{minipage}
\end{center}
\end{merksatz}
